\begin{figure}[ht]
\centering
\begin{tikzpicture}[x=0.32cm, y=0.32cm, font=\footnotesize]

% Axis
\draw[->, thick] (0,0) -- (38,0) node[right] {$\log_{10}(1/x)$};
\draw[->, thick] (0,0) -- (0,18) node[above, align=center] {decimal digits of\\correct result};

% Tick marks for x = 10^{-k}, k=0..16
\foreach \k in {0,2,4,6,8,10,12,14,16} {
  \pgfmathsetmacro{\xp}{2*\k}
  \draw (\xp,0) -- (\xp,-0.3);
  \node[anchor=north, font=\scriptsize] at (\xp,-0.3) {$10^{-\k}$};
}
% y ticks
\foreach \d in {0,4,8,12,16} {
  \draw (0,\d) -- (-0.3,\d);
  \node[anchor=east, font=\scriptsize] at (-0.3,\d) {\d};
}

% Reference: ideal 16 digits constant line
\draw[thick, engblue, dashed] (0,16) -- (32,16) node[right, font=\scriptsize, engblue] {Taylor $x + x^{3}/6$};

% Catastrophic cancellation: naive sin(x) - x + x^3/6 in double
% Digits lost ~ 2k for x = 10^{-k}; clamp at 0
% Sample points (k, digits) approximately: (0,16) (2,12) (4,8) (6,4) (7,2) (8,0)
\draw[thick, engred] plot[smooth] coordinates {
  (0,16) (2,14) (4,10) (6,6) (8,2) (10,0) (16,0) (32,0)
};
\node[engred, font=\scriptsize, anchor=west] at (10.5,0.8)
  {naive double-precision evaluation};

% Annotated region of cancellation
\draw[thick, engamber, fill=engamber!20, opacity=0.4]
  (10,0) rectangle (32,16);
\node[anchor=center, font=\scriptsize] at (21,10)
  {cancellation region: $\sin x \approx x$ to many digits;};
\node[anchor=center, font=\scriptsize] at (21,8.5)
  {the difference $\sin x - x$ in double};
\node[anchor=center, font=\scriptsize] at (21,7)
  {underflows the significand of $\sin x$.};

\end{tikzpicture}
\caption[Catastrophic cancellation in $\sin(x) - x + x^{3}/6$]{Catastrophic
cancellation in the naive evaluation of $\sin(x) - x + x^{3}/6$ at small
$x$ in IEEE 754 binary64. The Taylor-rewritten form preserves all 16
digits of precision; the naive form loses approximately two digits per
decade of $x$ because the subtraction $\sin(x) - x$ cancels the leading
identical digits of two nearly equal numbers. By $x \approx 10^{-10}$
the naive result has no correct digits. The remedy is rewriting before
evaluating: the cancellation lives in the algorithm's choice, not in the
arithmetic. Source: authors' computation; analytic prediction
$\sim 2 \log_{10}(1/x)$ digits lost.}
\label{fig:vol02:ch16:catastrophic-cancellation}
\end{figure}
